% from ega4-11-fr.tex
\oldpage[IV-3]{130}

Le $C_0$-module $\operatorname{gr}_{\mathfrak J_0}^*(M_0)$ est engendré par
un nombre fini d'éléments homogènes de degré 1, de sorte qu'il y a un
épimorphisme de $C_0$-modules gradués
\[
u_0:C_0^r\longrightarrow\operatorname{gr}_{\mathfrak J_0}^*(M_0).
\]
Comme $C=C_0\otimes_{A_0}A$ et que l'homomorphisme canonique
$\operatorname{gr}_{\mathfrak J_0}^*(M_0)\otimes_{A_0}A
\to\operatorname{gr}_{\mathfrak J}^*(M)$ est surjectif, on déduit donc de
$u_0$ un épimorphisme de $C$-modules gradués
\[
u:C^r\xrightarrow{u_0\otimes1_A}
\operatorname{gr}_{\mathfrak J_0}^*(M_0)\otimes_{A_0}A
\longrightarrow\operatorname{gr}_{\mathfrak J}^*(M),
\]
et tout revient à voir que le $C$-module gradué $Q=\ker(u)$ est \emph{de type
fini}.

\begin{lemma}[11.2.9.6]
\label{IV.11.2.9.6-fr}
Sous les hypothèses précédentes, soit $\mathfrak R_0$ un idéal de $A_0$~;
posons
\[
A'_0=A_0/\mathfrak R_0,\qquad
B'_0=B_0\otimes_{A_0}A'_0=B_0/\mathfrak R_0B_0,\qquad
\mathfrak J'_0=\mathfrak J_0B'_0,
\]
et $\mathfrak R=\mathfrak R_0A$, $A'=A/\mathfrak R$, et supposons en outre
que $\operatorname{gr}_{\mathfrak J'_0}^*(M_0\otimes_{A_0}A'_0)$ soit un
$A'_0$-module plat. Alors il existe un nombre fini d'éléments $q_h$
($1\leq h\leq s$) de $Q$ tels que, pour tout idéal premier
$\mathfrak p\supset\mathfrak RB$ de $B$, les images canoniques des $q_h$ dans
$Q_\mathfrak p$ engendrent le $C_\mathfrak p$-module $Q_\mathfrak p$.
\end{lemma}

Supposons ce lemme démontré, et notons qu'on peut encore l'appliquer en
remplaçant $A_0$ par $(A_0)_{t_0}$, $B_0$, $\mathfrak J_0$, $M_0$,
$A_\lambda$ et $A$ par les objets correspondants $(B_0)_{t_0}$,
$(\mathfrak J_0)_{t_0}$, $(M_0)_{t_0}$, $(A_\lambda)_{t_0}$ et $A_{t_0}$ pour
tout $t_0\in A_0$, puisque $(A_0)_{t_0}$ est un $A_0$-module plat. On
appliquera alors ce lemme en remplaçant successivement $A_0$ par chacun des
anneaux $(A_0)_{f_{ik}}$, $\mathfrak R_0$ étant remplacé par l'idéal
$\mathfrak R_{ik}$ définissant le sous-schéma fermé de $D(f_{ik})$ induit par
$S_{0i}$ sur l'ouvert $S_{0i}\cap D(f_{ik})$ de $S_{0i}$. Comme l'hypothèse
de platitude de (11.2.9.6) est vérifiée dans chacun de ces cas en raison de
(11.2.9.5), on obtient ainsi pour chaque couple $(i,k)$ une famille finie
$(a_{ikh})_{1\leq h\leq l}$ d'éléments de $Q_{f_{ik}}$ dont les images dans
$Q_{\mathfrak p_{ik}}$ engendrent ce $C_{\mathfrak p_{ik}}$-module, pour tout
idéal premier $\mathfrak p_{ik}$ de $B_{f_{ik}}$ contenant
$\mathfrak R_{ik}B_{f_{ik}}$. On peut écrire, pour un entier $d>0$ convenable,
$a_{ikh}=b_{ikh}/f_{ik}^d$ pour tous les indices $i$, $k$, $h$, avec
$b_{ikh}\in Q$. Puisque les $S_{0i}$ recouvrent $S_0$, tout idéal premier
$\mathfrak p$ de $B$ est tel que son image dans $S_0$ appartienne à un
ensemble $S_{0i}\cap D(f_{ik})$, donc l'image $\mathfrak p_{ik}$ de
$\mathfrak p$ dans $B_{f_{ik}}$ contient $\mathfrak R_{ik}B_{f_{ik}}$. Comme
les images des $b_{ikh}$ ($1\leq h\leq l$) dans
$Q_{\mathfrak p_{ik}}=Q_\mathfrak p$ engendrent ce $C_\mathfrak p$-module,
on en conclut que les $b_{ikh}$
($1\leq i\leq m$, $1\leq k\leq n$, $1\leq h\leq l$) engendrent le
$C$-module $Q$ (Bourbaki, \emph{Alg. comm.}, chap.~II, \S~3, n$^\circ$~3,
th.~1).

Reste à démontrer le lemme (11.2.9.6). Posons
\[
C'=C\otimes_AA'=C/\mathfrak RC,\qquad
B'=B\otimes_AA',\qquad \mathfrak J'=\mathfrak JB',\qquad
M'=M\otimes_AA'.
\]
Par hypothèse $\operatorname{gr}_{\mathfrak J}^*(M)$ est un $A$-module plat,
donc $\operatorname{gr}_{\mathfrak J'}^*(M')$ s'identifie à
$\operatorname{gr}_{\mathfrak J}^*(M)\otimes_AA'$ (11.2.9.2), et on a la
suite exacte ($\mathbf 0_{\mathrm I}$, 6.1.2)
\begin{equation}
0\longrightarrow Q\otimes_AA'\longrightarrow C'^{r}
\xrightarrow{u\otimes1_{A'}}\operatorname{gr}_{\mathfrak J'}^*(M')
\longrightarrow0.
\label{IV.11.2.9.7-fr}
\tag{11.2.9.7}
\end{equation}
Posons $C'_0=C_0\otimes_{A_0}A'_0=C_0/\mathfrak R_0C_0$ et considérons
l'épimorphisme de $C'_0$-modules gradués
\[
u'_0:C'_0{}^{r}\xrightarrow{u_0\otimes1_{A'_0}}
\operatorname{gr}_{\mathfrak J_0}^*(M_0)\otimes_{A_0}A'_0
\longrightarrow
\operatorname{gr}_{\mathfrak J'_0}^*(M_0\otimes_{A_0}A'_0).
\]
Soit $Q_0=\ker(u'_0)$~; utilisant le fait que
$\operatorname{gr}_{\mathfrak J'_0}^*(M_0\otimes_{A_0}A'_0)$ est un
$A'_0$-module plat, on voit que son produit tensoriel sur $A'_0$ par $A'$
s'identifie à $\operatorname{gr}_{\mathfrak J'}^*(M')$ (11.2.9.2) et on a la
suite exacte ($\mathbf 0_{\mathrm I}$, 6.1.2)
\begin{equation}
0\longrightarrow Q_0\otimes_{A'_0}A'\longrightarrow C'^{r}
\xrightarrow{u\otimes1_{A'}}\operatorname{gr}_{\mathfrak J'}^*(M')
\longrightarrow0.
\label{IV.11.2.9.8-fr}
\tag{11.2.9.8}
\end{equation}

